\uuid{CzJB}
\titre{Changement de l'inconnue}
\auteur{hanani ; mzari}
\organisation{exo7}
\licence{CC BY-NC-SA 4.0}
\datecreate{2020-01-20}
\niveau{L1}
\module{Analyse}
\chapitre{Equation différentielle}
\sousChapitre{Résolution d'équation différentielle du deuxième ordre}
\difficulte{3}
\theme{Equations du second ordre, changement de l'inconnue}
\exo7qcmid{lille-2.equations-differentielles.033}

\contenu{
\begin{question}
\qclassification{225.02}
\qtags{motcle=changement de l'inconnue}
On considère les équations différentielles
$$(E_1)\; :\; xy''+2(x+1)y'+(x+2)y=2\mathrm{e}^{-x}\quad\mbox{et}\quad (E_2)\; :\; y''+2y'+y=2\mathrm{e}^{-x}.$$
Parmi les affirmations suivantes, cocher celles qui sont vraies :
\begin{answers}
\good{Si $y$ est une solution de $(E_1)$, alors $\displaystyle z=xy$ est une solution de $(E_2)$.}
\bad{$(E_2)$ admet une solution particulière sous la forme $\displaystyle y_0=a\mathrm{e}^{-x}$, $a\in \Rr$.}
\good{La solution générale de $(E_1)$ sur $]0,+\infty[$ est $\displaystyle y=\left(x+a+\frac{b}{x}\right)\mathrm{e}^{-x}$, $a,b\in \Rr$.}
\bad{Toute solution de $(E_1)$ sur $]0,+\infty[$ se prolonge par continuité en $0$.}
\end{answers}
\begin{explanations}
On a : $z'=xy'+y$, $z''=xy''+2y'$ et $z''+2z'+z=xy''+2(x+1)y'+(x+2)y$. La solution générale de l'équation homogène associée à $(E_2)$ est : $Z=(ax+b)\mathrm{e}^{-x}$ et $z_0=x^2\mathrm{e}^{-x}$ est une solution particulière de $(E_2)$. Donc la solution générale de $(E_2)$ est $z=Z+z_0$. Ainsi, la solution générale de $(E_1)$ sur $]0,+\infty[$ est : $\displaystyle y=\left(x+a+\frac{b}{x}\right)\mathrm{e}^{-x}$, $a,b\in \Rr$. Une telle solution se prolonge par continuité en $0$ si et seulement si $b=0$.
\end{explanations}
\end{question}
}
